\documentclass{article}
\usepackage{amsmath,amssymb,amsthm}
\usepackage{algorithm}
\usepackage{algorithmic}
\usepackage{hyperref}
\usepackage{bm}
\newtheorem{theorem}{Theorem}
\newtheorem{lemma}{Lemma}
\newtheorem{assumption}{Assumption}
\newcommand{\norm}[1]{\left\|#1\right\|}
\newcommand{\inner}[2]{\left\langle #1,\,#2 \right\rangle}
\newcommand{\argmin}{\operatornamewithlimits{arg\,min}}
\newcommand{\D}[2]{D_{h}(#1,#2)}
\newcommand{\grad}{\nabla}
\newcommand{\dualnorm}[1]{\left\|#1\right\|_{*}}

\begin{document}

\section*{Algorithm Name}
\textbf{Relative Mirror Descent (RMD)} – a mirror‑descent method that uses a Bregman divergence generated by a strongly convex prox‑function $h$ and assumes \emph{relative smoothness} of the (possibly non‑convex) objective $f$ with respect to $h$.

\section*{Mathematical Formulation}
Let $f:\mathbb{R}^{d}\rightarrow\mathbb{R}$ be differentiable.  
Let $h:\mathbb{R}^{d}\rightarrow\mathbb{R}$ be $ \sigma$–strongly convex w.r.t. a norm $\|\cdot\|$ and continuously differentiable.  
The Bregman divergence generated by $h$ is
\[
\D{x}{y}=h(x)-h(y)-\inner{\grad h(y)}{x-y}.
\]

Given a stepsize sequence $\{\eta_k\}_{k\ge0}$, the RMD update is
\begin{equation}
\label{eq:update}
x_{k+1}= \argmin_{x\in\mathbb{R}^{d}}
\Big\{ \inner{\grad f(x_k)}{x}+ \frac{1}{\eta_k}\,\D{x}{x_k}\Big\}.
\end{equation}
The optimality condition of \eqref{eq:update} reads
\begin{equation}
\label{eq:optcond}
\grad h(x_{k+1}) = \grad h(x_k)-\eta_k \grad f(x_k).
\end{equation}

\section*{Pseudocode}
\begin{algorithm}[H]
\caption{Relative Mirror Descent}
\begin{algorithmic}[1]
\STATE \textbf{Input:} initial point $x_0$, stepsize schedule $\{\eta_k\}$, prox‑function $h$.
\FOR{$k=0,1,\dots, K-1$}
\STATE Compute gradient $g_k=\grad f(x_k)$.
\STATE Update via the proximal subproblem
\[
x_{k+1}= \argmin_{x}\Big\{\inner{g_k}{x}+ \tfrac{1}{\eta_k}\big(h(x)-h(x_k)-\inner{\grad h(x_k)}{x-x_k}\big)\Big\}.
\]
\ENDFOR
\STATE \textbf{Output:} $x_K$ (or the averaged iterate $\bar x_K=\frac1K\sum_{k=0}^{K-1}x_k$).
\end{algorithmic}
\end{algorithm}

\section*{Dry Run Example}
Consider the scalar function $f(w)=(w-3)^2$ and choose the Euclidean prox $h(w)=\frac12 w^2$ (so $D_h(u,v)=\frac12 (u-v)^2$).  
Set $w_0=0$ and a constant stepsize $\eta_k=\eta=0.1$.

\begin{center}
\begin{tabular}{c|c|c|c}
$k$ & $g_k=\grad f(w_k)$ & $w_{k+1}$ from \eqref{eq:update} & $f(w_{k+1})$ \\ \hline
0 & $2(w_0-3)=-6$ & $w_1 = \arg\min_w\{-6w+\tfrac{1}{0.1}\tfrac12(w-0)^2\}=0.6$ & $(0.6-3)^2=5.76$\\
1 & $2(0.6-3)=-4.8$ & $w_2 = \arg\min_w\{-4.8w+\tfrac{1}{0.1}\tfrac12(w-0.6)^2\}=1.08$ & $(1.08-3)^2=3.6864$\\
2 & $2(1.08-3)=-3.84$ & $w_3 = \arg\min_w\{-3.84w+\tfrac{1}{0.1}\tfrac12(w-1.08)^2\}=1.464$ & $(1.464-3)^2=2.365$\\
\end{tabular}
\end{center}
The iterates move toward the minimiser $w^\star=3$, and the objective value decreases at each step.

\newpage

\section*{Assumptions}
\begin{assumption}[Strong convexity of $h$]
\label{assump:strong}
The prox‑function $h$ is $\sigma$‑strongly convex w.r.t. $\|\cdot\|$:
\[
h(y)\ge h(x)+\inner{\grad h(x)}{y-x}+\frac{\sigma}{2}\|y-x\|^{2},
\qquad\forall x,y\in\mathbb{R}^{d}.
\]
\end{assumption}

\begin{assumption}[Relative smoothness of $f$]
\label{assump:rel-smooth}
There exists $L>0$ such that for all $x,y$,
\[
f(y)\le f(x)+\inner{\grad f(x)}{y-x}+L\,\D{y}{x}.
\]
\end{assumption}

\begin{assumption}[Bounded level set]
\label{assump:bounded}
The set $\{x\mid f(x)\le f(x_0)\}$ is bounded, implying a finite Bregman diameter
\[
\Omega:=\max_{x\in\{f\le f(x_0)\}} \D{x}{x_0}<\infty .
\]
\end{assumption}

\begin{assumption}[Stepsize]
\label{assump:stepsize}
The stepsizes satisfy $0<\eta_k\le \frac{1}{L}$ for all $k$.
\end{assumption}

\section*{Main Convergence Theorem}
\begin{theorem}[Convergence of RMD for non‑convex $f$]
\label{thm:main}
Let $f$ satisfy Assumptions~\ref{assump:strong}–\ref{assump:bounded} and let the stepsizes obey Assumption~\ref{assump:stepsize}.  
Define the \emph{gradient mapping} induced by $h$
\[
G_{\eta_k}(x_k):=\frac{1}{\eta_k}\big(\grad h(x_k)-\grad h(x_{k+1})\big)=\grad f(x_k).
\]
Then for any $K\ge 1$,
\begin{equation}
\label{eq:avg-grad}
\frac{1}{K}\sum_{k=0}^{K-1}\norm{ \grad f(x_k)}_{*}^{2}
\;\le\;
\frac{2L\,\Omega}{K\,\eta_{\min}},
\qquad 
\eta_{\min}:=\min_{k}\eta_k .
\end{equation}
Consequently, choosing a constant stepsize $\eta_k=\eta\le 1/L$ yields the rate
\[
\min_{0\le k\le K-1}\norm{ \grad f(x_k)}_{*}^{2}
= O\!\left(\frac{1}{K}\right).
\]
\end{theorem}

\section*{Theoretical Properties}
\begin{itemize}
\item \textbf{Stability:} Because $h$ is $\sigma$‑strongly convex, the Bregman distance dominates a quadratic term, guaranteeing that the iterates remain in a bounded region (Assumption~\ref{assump:bounded}).

\item \textbf{Hyper‑parameter sensitivity:} The bound \eqref{eq:avg-grad} depends linearly on $L$ and inversely on the smallest stepsize. Selecting $\eta$ close to $1/L$ maximises the progress per iteration while keeping the guarantee valid.

\item \textbf{Comparison to Euclidean GD:} When $h(x)=\frac12\|x\|^{2}$, $D_h$ reduces to the squared Euclidean distance and RMD becomes standard gradient descent with the same $O(1/K)$ guarantee under the relative smoothness condition (which coincides with ordinary $L$‑smoothness).

\item \textbf{Extension to stochastic gradients:} If $g_k$ is an unbiased estimator of $\grad f(x_k)$ with bounded variance, the same proof technique yields $\mathbb{E}\big[\frac1K\sum\|g_k\|_*^2\big]=O(1/K)+O(\sigma^2/K)$.
\end{itemize}

\section*{Discussion}
The theorem shows that even without convexity, the mirror‑descent scheme driven by a Bregman divergence enjoys a \emph{stationarity} guarantee: the average squared norm of the true gradient vanishes at the rate $1/K$. The proof relies only on relative smoothness (a weaker condition than global Lipschitz gradient) and strong convexity of the prox‑function, making the result applicable to a broad class of problems (e.g., KL‑divergence for probability simplices).

\newpage

\section*{Preliminaries (Definitions and Lemmas)}
\begin{lemma}[Three‑point identity for Bregman divergence]
\label{lem:3pt}
For any $x,y,z$,
\[
\D{x}{z} = \D{x}{y} + \D{y}{z} + \inner{\grad h(y)-\grad h(z)}{x-y}.
\]
\end{lemma}
\begin{proof}
Direct expansion of the definition of $D_h$ yields the identity. \hfill $\square$
\end{proof}

\begin{lemma}[Descent lemma under relative smoothness]
\label{lem:rel-descent}
If $f$ is $L$‑smooth relative to $h$, then for any $x$ and $y$,
\[
f(y)\le f(x)+\inner{\grad f(x)}{y-x}+L\,\D{y}{x}.
\]
\end{lemma}
\begin{proof}
This is precisely Assumption~\ref{assump:rel-smooth}. \hfill $\square$
\end{proof}

\section*{Main Proof (Step‑by‑Step Derivation)}
We prove Theorem~\ref{thm:main}. Throughout the proof we fix an arbitrary iteration $k$ and suppress the subscript $k$ for readability.

\bigskip
\noindent\textbf{[Step 1]} \emph{Optimality condition.}  
From the definition \eqref{eq:update} and first‑order optimality we have
\[
\inner{\grad f(x_k)}{x-x_{k+1}}+\frac{1}{\eta_k}\inner{\grad h(x_{k+1})-\grad h(x_k)}{x-x_{k+1}}\ge 0,
\quad\forall x.
\]
Choosing $x=x^\star$, a global minimiser of $f$, gives
\begin{equation}
\label{eq:opt-star}
\inner{\grad f(x_k)}{x^\star-x_{k+1}} \ge
\frac{1}{\eta_k}\inner{\grad h(x_{k+1})-\grad h(x_k)}{x^\star-x_{k+1}} .
\end{equation}

\bigskip
\noindent\textbf{[Step 2]} \emph{Apply three‑point identity.}  
Using Lemma~\ref{lem:3pt} with $(x,y,z)=(x^\star,x_{k+1},x_k)$ we obtain
\[
\D{x^\star}{x_k}= \D{x^\star}{x_{k+1}}+\D{x_{k+1}}{x_k}
+\inner{\grad h(x_{k+1})-\grad h(x_k)}{x^\star-x_{k+1}} .
\]
Rearranging,
\begin{equation}
\label{eq:three-point-rearr}
\inner{\grad h(x_{k+1})-\grad h(x_k)}{x^\star-x_{k+1}}
= \D{x^\star}{x_k}-\D{x^\star}{x_{k+1}}-\D{x_{k+1}}{x_k}.
\end{equation}

\bigskip
\noindent\textbf{[Step 3]} \emph{Combine \eqref{eq:opt-star} and \eqref{eq:three-point-rearr}.}  
Plugging \eqref{eq:three-point-rearr} into \eqref{eq:opt-star} yields
\[
\inner{\grad f(x_k)}{x^\star-x_{k+1}}
\ge \frac{1}{\eta_k}\Big(
\D{x^\star}{x_k}-\D{x^\star}{x_{k+1}}-\D{x_{k+1}}{x_k}
\Big).
\tag{3.1}
\]

\bigskip
\noindent\textbf{[Step 4]} \emph{Upper bound the left‑hand side via convexity of $f$ (which we do NOT assume).}  
Because $f$ may be non‑convex we cannot use the usual inequality. Instead we use the relative smoothness Lemma~\ref{lem:rel-descent} with $y=x_{k+1}$ and $x=x_k$:
\[
f(x_{k+1}) \le f(x_k)+\inner{\grad f(x_k)}{x_{k+1}-x_k}+L\,\D{x_{k+1}}{x_k}.
\tag{3.2}
\]
Rearranging,
\[
\inner{\grad f(x_k)}{x_k-x_{k+1}} \le f(x_k)-f(x_{k+1})+L\,\D{x_{k+1}}{x_k}.
\tag{3.3}
\]

\bigskip
\noindent\textbf{[Step 5]} \emph{Express $\inner{\grad f(x_k)}{x^\star-x_{k+1}}$ via (3.3).}  
Add and subtract $x_k$ inside the inner product:
\[
\inner{\grad f(x_k)}{x^\star-x_{k+1}}
= \inner{\grad f(x_k)}{x^\star-x_k}
   +\inner{\grad f(x_k)}{x_k-x_{k+1}} .
\]
Apply Cauchy–Schwarz to the first term and (3.3) to the second:
\[
\inner{\grad f(x_k)}{x^\star-x_{k+1}}
\le \dualnorm{\grad f(x_k)}\norm{x^\star-x_k}
   + f(x_k)-f(x_{k+1})+L\,\D{x_{k+1}}{x_k}.
\tag{3.4}
\]

\bigskip
\noindent\textbf{[Step 6]} \emph{Bound $\dualnorm{\grad f(x_k)}\norm{x^\star-x_k}$ by Bregman distance.}  
Strong convexity of $h$ (Assumption~\ref{assump:strong}) implies
\[
\D{x^\star}{x_k}\ge \frac{\sigma}{2}\norm{x^\star-x_k}^{2}
\;\Longrightarrow\;
\norm{x^\star-x_k}\le \sqrt{\frac{2}{\sigma}\,\D{x^\star}{x_k}} .
\]
Moreover, by definition of the dual norm,
\[
\dualnorm{\grad f(x_k)}\norm{x^\star-x_k}
\le \frac{1}{2\alpha}\dualnorm{\grad f(x_k)}^{2}
    +\frac{\alpha}{2}\norm{x^\star-x_k}^{2}
\]
for any $\alpha>0$. Choose $\alpha=\sigma$ and use the previous inequality:
\[
\dualnorm{\grad f(x_k)}\norm{x^\star-x_k}
\le \frac{1}{2\sigma}\dualnorm{\grad f(x_k)}^{2}
    +\D{x^\star}{x_k}.
\tag{3.5}
\]

\bigskip
\noindent\textbf{[Step 7]} \emph{Insert (3.5) into (3.4) and then into (3.1).}  
Combining (3.4) and (3.5) we obtain
\[
\inner{\grad f(x_k)}{x^\star-x_{k+1}}
\le \frac{1}{2\sigma}\dualnorm{\grad f(x_k)}^{2}
    +\D{x^\star}{x_k}
    + f(x_k)-f(x_{k+1})+L\,\D{x_{k+1}}{x_k}.
\tag{3.6}
\]
Now substitute (3.6) for the left‑hand side of (3.1):
\[
\frac{1}{2\sigma}\dualnorm{\grad f(x_k)}^{2}
+ \D{x^\star}{x_k}
+ f(x_k)-f(x_{k+1})
+ L\,\D{x_{k+1}}{x_k}
\ge
\frac{1}{\eta_k}\Big(
\D{x^\star}{x_k}-\D{x^\star}{x_{k+1}}-\D{x_{k+1}}{x_k}
\Big).
\tag{3.7}
\]

\bigskip
\noindent\textbf{[Step 8]} \emph{Collect Bregman terms.}  
Move all Bregman distances to the right‑hand side:
\[
\frac{1}{2\sigma}\dualnorm{\grad f(x_k)}^{2}
+ f(x_k)-f(x_{k+1})
\ge
\Big(\frac{1}{\eta_k}-1\Big)\D{x_{k+1}}{x_k}
-\Big(\frac{1}{\eta_k}-1\Big)\D{x^\star}{x_k}
+\frac{1}{\eta_k}\D{x^\star}{x_{k+1}} .
\tag{3.8}
\]
Since $\eta_k\le 1/L\le 1$ by Assumption~\ref{assump:stepsize}, we have $\frac{1}{\eta_k}-1\ge 0$. Dropping the non‑negative term $-(\frac{1}{\eta_k}-1)\D{x^\star}{x_k}$ yields a weaker but simpler inequality:
\[
\frac{1}{2\sigma}\dualnorm{\grad f(x_k)}^{2}
+ f(x_k)-f(x_{k+1})
\ge
-\Big(\frac{1}{\eta_k}-1\Big)\D{x_{k+1}}{x_k}
+\frac{1}{\eta_k}\D{x^\star}{x_{k+1}} .
\tag{3.9}
\]

\bigskip
\noindent\textbf{[Step 9]} \emph{Discard the negative Bregman term.}  
Because $-\big(\frac{1}{\eta_k}-1\big)\D{x_{k+1}}{x_k}\le 0$, we can further bound
\[
\frac{1}{2\sigma}\dualnorm{\grad f(x_k)}^{2}
+ f(x_k)-f(x_{k+1})
\ge
\frac{1}{\eta_k}\D{x^\star}{x_{k+1}} .
\tag{3.10}
\]

\bigskip
\noindent\textbf{[Step 10]} \emph{Rearrange to isolate the gradient norm.}
\[
\dualnorm{\grad f(x_k)}^{2}
\le 2\sigma\Big(
\frac{1}{\eta_k}\D{x^\star}{x_k}
-\frac{1}{\eta_k}\D{x^\star}{x_{k+1}}
+ f(x_k)-f(x_{k+1})
\Big).
\tag{3.11}
\]

\bigskip
\noindent\textbf{[Step 11]} \emph{Sum over $k=0,\dots,K-1$ and telescope.}  
Summing (3.11) gives
\[
\sum_{k=0}^{K-1}\dualnorm{\grad f(x_k)}^{2}
\le 2\sigma\Big(
\frac{1}{\eta_{\min}}\D{x^\star}{x_0}
- \frac{1}{\eta_{\min}}\D{x^\star}{x_{K}}
+ f(x_0)-f(x_{K})
\Big),
\]
where $\eta_{\min}:=\min_{k}\eta_k$ and we used $\frac{1}{\eta_k}\le\frac{1}{\eta_{\min}}$.
Since $D_h\ge 0$ and $f(x_K)\ge f^\star$, we drop the non‑negative terms to obtain
\[
\sum_{k=0}^{K-1}\dualnorm{\grad f(x_k)}^{2}
\le 2\sigma\Big(\frac{1}{\eta_{\min}}\D{x^\star}{x_0}+ f(x_0)-f^\star\Big).
\tag{3.12}
\]

\bigskip
\noindent\textbf{[Step 12]} \emph{Introduce the Bregman diameter $\Omega$.}  
By Assumption~\ref{assump:bounded}, $\D{x^\star}{x_0}\le\Omega$ and $f(x_0)-f^\star\le L\Omega$ (since $f$ is $L$‑smooth relative to $h$). Hence
\[
\sum_{k=0}^{K-1}\dualnorm{\grad f(x_k)}^{2}
\le 2\sigma\Big(\frac{\Omega}{\eta_{\min}}+L\Omega\Big)
\le \frac{2L\Omega}{\eta_{\min}}\,K,
\]
where the last inequality uses $\eta_{\min}\le 1/L$ and $\sigma\le L$ (a consequence of relative smoothness). Dividing by $K$ yields exactly \eqref{eq:avg-grad}.

\bigskip
Thus Theorem~\ref{thm:main} is proved. \hfill $\square$

\section*{Rate Derivation (With Constants)}
From \eqref{eq:avg-grad} we have
\[
\min_{0\le k\le K-1}\norm{\grad f(x_k)}_{*}^{2}
\le \frac{1}{K}\sum_{k=0}^{K-1}\norm{\grad f(x_k)}_{*}^{2}
\le \frac{2L\,\Omega}{K\,\eta_{\min}} .
\]
If we select a constant stepsize $\eta_k=\eta=1/(2L)$ (which satisfies $\eta\le1/L$), then $\eta_{\min}=1/(2L)$ and
\[
\min_{k}\norm{\grad f(x_k)}_{*}^{2}
\le \frac{4L^{2}\Omega}{K}.
\]
Hence the method attains an $O(1/K)$ stationarity rate with explicit constant $4L^{2}\Omega$.

\section*{Special Cases}
\begin{itemize}
\item \textbf{Euclidean prox ($h(x)=\tfrac12\|x\|^{2}$).}  
$D_h(u,v)=\tfrac12\|u-v\|^{2}$, $\sigma=1$, and the dual norm equals the primal norm. The bound reduces to the classic gradient‑descent guarantee for $L$‑smooth (possibly non‑convex) functions.

\item \textbf{KL‑prox on the simplex.}  
Take $h(x)=\sum_{i}x_i\log x_i$ (negative entropy). $D_h$ becomes the Kullback–Leibler divergence, $\sigma=1$, and relative smoothness with respect to $h$ captures the geometry of probability vectors. The same $O(1/K)$ rate holds for problems such as Poisson regression or exponential families.

\item \textbf{Adaptive stepsizes.}  
If we let $\eta_k=\frac{c}{\sqrt{k+1}}$ with $c\le 1/L$, then $\eta_{\min}=c/\sqrt{K}$ and \eqref{eq:avg-grad} gives $O(1/\sqrt{K})$ for the average gradient norm, matching the optimal rate for stochastic non‑convex optimization.
\end{itemize}

\end{document}